\documentclass[12pt]{article}
\usepackage{amsmath}
\usepackage{amsfonts}
\usepackage{amssymb}
\usepackage{geometry}
\usepackage{physics}
\usepackage{tikz}
\geometry{a4paper, margin=0.7in}


\title{02YMECH - Solution 4}
\date{Assigned for the week of Oct 13, 2025}
\author{}
\begin{document}
\maketitle

\section*{Solutions}

\begin{enumerate}

% 1
\item \textbf{Time to first maximum displacement (undamped oscillator).}

The motion is
\[
x(t) = A \cos(\omega_0 t + \phi),
\]
where $A$, $\omega_0$, and $\phi$ are constants. The velocity is
\[
v(t) = \dot{x}(t) = -A \omega_0 \sin(\omega_0 t + \phi).
\]

Maximum displacement in simple harmonic motion occurs when $|x(t)| = A$, i.e. when the cosine is $\pm 1$. At those instants, the velocity is zero.

We are told the initial position and velocity:
\[
x(0) = x_0 = A \cos \phi,
\qquad
v(0) = v_0 = -A \omega_0 \sin \phi.
\]

From these,
\[
\cos \phi = \frac{x_0}{A},
\qquad
\sin \phi = - \frac{v_0}{A \omega_0}.
\]

We want the first time $t>0$ such that $x(t)$ reaches an \emph{extreme} (turning point), i.e. when $v(t)=0$ \emph{and} $x$ is a maximum in magnitude.

The condition $v(t)=0$ gives
\[
\sin(\omega_0 t + \phi) = 0
\quad \Longrightarrow \quad
\omega_0 t + \phi = n\pi,
\quad n \in \mathbb{Z}.
\]

Thus the candidate times are
\[
t_n = \frac{n\pi - \phi}{\omega_0}.
\]

At those instants,
\[
x(t_n) = A \cos(n\pi) = A(-1)^n.
\]
So for even $n$ ($n=0,2,4,\dots$), $x = +A$; for odd $n$ ($n=1,3,5,\dots$), $x = -A$. Both $\pm A$ are extrema of equal magnitude.

We want the \emph{first} $t>0$. The smallest $t_n>0$ is obtained by choosing the smallest integer $n$ such that
\[
t_n = \frac{n\pi - \phi}{\omega_0} > 0.
\]

This gives the rule:
\[
t_{\text{first}} = \min_{n \in \mathbb{Z}} \left\{ \frac{n\pi - \phi}{\omega_0} \, \Big| \, \frac{n\pi - \phi}{\omega_0} > 0 \right\}.
\]

In words: find the next zero of the sine after $t=0$. Equivalently, if $\phi \in (-\pi,\pi]$,
\[
t_{\text{first}} =
\begin{cases}
\displaystyle \frac{-\phi}{\omega_0}, & \text{if } \phi < 0, \\[1em]
\displaystyle \frac{\pi - \phi}{\omega_0}, & \text{if } \phi \ge 0~.
\end{cases}
\]

We can also express $\phi$ in terms of $x_0,v_0$ using
\[
\phi = \arctan2\!\left(-\frac{v_0}{A\omega_0}, \frac{x_0}{A}\right),
\]
where $\arctan2(y,x)$ is the correct quadrant angle with $\cos\phi = x_0/A$ and $\sin\phi = -v_0/(A\omega_0)$.

Summary: the first time the oscillator reaches an endpoint ($|x|=A$) after $t=0$ is given by the smallest positive member of $t_n = (n\pi-\phi)/\omega_0$.

Physically: the oscillator reaches maximum displacement whenever its velocity first becomes zero.


% 2
\item \textbf{Average kinetic and potential energy over one cycle.}

For a harmonic oscillator of mass $m$ and angular frequency $\omega_0$, we can write
\[
x(t) = A \cos(\omega_0 t + \phi).
\]
Then
\[
v(t) = \dot{x}(t) = -A \omega_0 \sin(\omega_0 t + \phi).
\]

The kinetic energy is
\[
K(t) = \frac{1}{2} m v^2(t)
= \frac{1}{2} m A^2 \omega_0^2 \sin^2(\omega_0 t + \phi).
\]

The potential energy (spring potential) is
\[
U(t) = \frac{1}{2} k x^2(t)
= \frac{1}{2} k A^2 \cos^2(\omega_0 t + \phi),
\]
with $k = m \omega_0^2$ for a simple harmonic oscillator. Thus
\[
U(t) = \frac{1}{2} m \omega_0^2 A^2 \cos^2(\omega_0 t + \phi).
\]

The total mechanical energy is constant:
\[
E = K(t)+U(t)
= \frac{1}{2} m \omega_0^2 A^2.
\]

We now compute time averages over one full period $T = \frac{2\pi}{\omega_0}$:
\[
\langle K \rangle
= \frac{1}{T} \int_0^T K(t)\, dt,
\qquad
\langle U \rangle
= \frac{1}{T} \int_0^T U(t)\, dt.
\]

We use the facts that over any full period,
\[
\langle \sin^2(\omega_0 t + \phi) \rangle = \frac{1}{2},
\qquad
\langle \cos^2(\omega_0 t + \phi) \rangle = \frac{1}{2}.
\]

Therefore,
\[
\langle K \rangle
= \frac{1}{2} m A^2 \omega_0^2 \cdot \frac{1}{2}
= \frac{1}{4} m A^2 \omega_0^2,
\]
\[
\langle U \rangle
= \frac{1}{2} m A^2 \omega_0^2 \cdot \frac{1}{2}
= \frac{1}{4} m A^2 \omega_0^2.
\]

But $E = \tfrac{1}{2} m \omega_0^2 A^2$, so
\[
\langle K \rangle = \frac{E}{2},
\qquad
\langle U \rangle = \frac{E}{2}.
\]

Conclusion: over one complete cycle, the oscillator spends half its energy in kinetic form (on average) and half in potential form. The averages are equal and each is $E/2$.


% 3
\item \textbf{Pendulum with a horizontally accelerating support.}

We have a pendulum of length $L$ and bob mass $m$. Initially it hangs straight down. Then the support point begins accelerating horizontally with constant acceleration $a$ (say to the right). After transients, the pendulum reaches a new equilibrium at some constant angle $\theta$ from the vertical, and then small oscillations about that equilibrium are possible.

We will work in the non-inertial frame accelerating with the support. In that frame, there is an effective ``pseudo-force'' of magnitude $ma$ acting on the bob opposite the acceleration of the support (i.e. to the left if the support accelerates to the right). Gravity $mg$ acts downward.

\begin{enumerate}
    \item \textit{New equilibrium angle $\theta$.}

    In the accelerating frame, equilibrium means the tension $\vec{T}$ in the string balances the vector sum of gravity and the horizontal pseudo-force. Let $\theta$ be the angle the string makes with the vertical, displaced in the \emph{opposite} direction of the pseudo-force so that the bob is ``behind'' the accelerating support.

    Force balance requires that the string aligns with the resultant of $mg$ downward and $ma$ horizontally. Geometrically,
    \[
    \tan \theta = \frac{\text{horizontal pseudo-force}}{\text{gravity}} = \frac{ma}{mg} = \frac{a}{g}.
    \]

    Thus
    \[
    \theta = \arctan\!\left(\frac{a}{g}\right).
    \]

    \item \textit{Period of small oscillations about this equilibrium.}

    For small oscillations about the new equilibrium, the pendulum behaves like a simple pendulum of length $L$ in a uniform effective gravitational field $\vec{g}_\text{eff}$ equal to the vector sum of the real gravity and the pseudo-gravity from acceleration.

    The magnitude of this effective gravitational field is
    \[
    g_\text{eff} = \sqrt{g^2 + a^2}.
    \]

    The restoring torque for small angular displacements $\delta\theta$ about equilibrium is proportional to $mg_\text{eff} L \, \delta\theta$, so the small-angle equation of motion is that of a simple pendulum with $g$ replaced by $g_\text{eff}$.

    Therefore, the small-oscillation angular frequency is
    \[
    \omega_{\text{small}} = \sqrt{\frac{g_\text{eff}}{L}}
    = \sqrt{\frac{\sqrt{g^2 + a^2}}{L}}.
    \]

    The corresponding period is
    \[
    T = \frac{2\pi}{\omega_{\text{small}}}
    = 2\pi \sqrt{\frac{L}{g_\text{eff}}}
    = 2\pi \sqrt{\frac{L}{\sqrt{g^2 + a^2}}}.
    \]

    Check limiting cases:
    \begin{itemize}
        \item If $a=0$, then $g_\text{eff}=g$ and $T=2\pi \sqrt{L/g}$, which is the usual pendulum period.
        \item If $a$ is large, the effective gravity increases, so the period decreases.
    \end{itemize}
\end{enumerate}


% 4
\item \textbf{Damped harmonic oscillator amplitude decay.}

The position of a damped oscillator is
\[
x(t) = A e^{-\beta t} \cos(\omega_1 t + \phi),
\]
where $A$ is the initial amplitude factor, $\beta>0$ is the damping coefficient (with units of 1/time), and $\omega_1$ is the damped angular frequency.

The amplitude envelope of the oscillation is
\[
A_{\text{env}}(t) = A e^{-\beta t}.
\]

We are asked: how long until the amplitude decreases to half its initial value?

We want $t = t_{1/2}$ such that
\[
A_{\text{env}}(t_{1/2}) = \frac{A}{2}.
\]

That is,
\[
A e^{-\beta t_{1/2}} = \frac{A}{2}
\quad\Longrightarrow\quad
e^{-\beta t_{1/2}} = \frac{1}{2}.
\]

Take the natural logarithm:
\[
-\beta t_{1/2} = \ln\!\left(\frac{1}{2}\right) = -\ln 2,
\]
so
\[
t_{1/2} = \frac{\ln 2}{\beta}.
\]

This $t_{1/2}$ is the amplitude ``half-life'' for the damped oscillator. Note that it depends only on $\beta$, not on $\omega_1$ or $A$.
\end{enumerate}


\end{document}
